\begin{helper}[Product Model Lipschitz Estimate]
Fix \(n,M\in\mathbb N\), \(I\subseteq\mathrm{Fin}(n)\),
\(\varepsilon,p_0\in\mathbb R\), an embedding
\[
\pi:\mathrm{Fin}(n)\hookrightarrow\mathrm{Fin}(M)\times\mathrm{Fin}(M),
\]
and set
\[
c_{\mathrm{prod}}=\frac12 n^{4\varepsilon}.
\]
If two product assignments \(v\) and \(v'\) agree at every coordinate except
possibly one, then
\[
\left|
\noderef{FixedSetProductModelVar}(I,\varepsilon,p_0,\pi,v)
-
\noderef{FixedSetProductModelVar}(I,\varepsilon,p_0,\pi,v')
\right|
\le c_{\mathrm{prod}}.
\]
\end{helper}
\begin{proof}
Unfold \(\noderef{FixedSetProductModelVar}\).  Each product coordinate
corresponds to one base vertex: if \(i<M\), it controls the first component of
the red external-neighborhood data at the red vertex \(x\); if \(M\le i<2M\),
it controls the second component of the blue external-neighborhood data at the
blue vertex \(x\).  Let \(x_i\) denote this one possibly changed base vertex.
For every other base vertex \(y\ne x_i\), the reconstructed graph adjacency
between \(y\) and the embedded projections of \(I\) is identical for \(v\) and
\(v'\).  Hence
\[
\noderef{TwoBiteX}(\omega_v,I,y)
=\noderef{TwoBiteX}(\omega_{v'},I,y),
\]
and the small-class predicate and final-pair contribution at \(y\) are also
identical.

For each assignment \(w\), let
\[
F_y(w)=
\noderef{TwoBiteFinalPairs}(\noderef{TwoBiteX}(\omega_w,I,y))
\]
after applying the distinct-red-and-blue-projection filter, and set
\(F_y(w)=\varnothing\) if \(y\) is not in the external small-class set
\[
S_I^{\mathrm{ext}}(\omega_w)
=\{z:\noderef{TwoBiteSmallClass}(\omega_w,\varepsilon,I,z)\}\setminus P_I.
\]
This convention includes the embedded-coordinate cases \(y\in P_R\) or
\(y\in P_B\): such \(y\) is removed by \(P_I\), so its contribution is empty
for both assignments.  With this notation the counted set for \(w\) is
\[
A(w)=\bigcup_y F_y(w).
\]
All terms with \(y\ne x_i\) are the same for \(v\) and \(v'\).  Let
\[
U=\bigcup_{y\ne x_i}F_y(v)=\bigcup_{y\ne x_i}F_y(v').
\]
Then
\[
A(v)=U\cup F_{x_i}(v),\qquad
A(v')=U\cup F_{x_i}(v').
\]
Since \(U\) is common,
\[
|A(v)|=|U|+|F_{x_i}(v)\setminus U|,\qquad
|A(v')|=|U|+|F_{x_i}(v')\setminus U|.
\]
Therefore
\[
\bigl||A(v)|-|A(v')|\bigr|
\le
\max\{|F_{x_i}(v)\setminus U|,\ |F_{x_i}(v')\setminus U|\}
\le
\max\{|F_{x_i}(v)|,\ |F_{x_i}(v')|\}.
\]
This is the exact set-cardinality comparison: only the one changed
coordinate's contribution can alter the union cardinality; the common union
does not create a second uncontrolled error term.

It remains to bound each possible changed contribution.  If \(x_i\) is not in
the external small class for the relevant assignment, then
\(F_{x_i}=\varnothing\).  If \(x_i\) is in the external small class, the
definition of \(\noderef{TwoBiteSmallClass}\) gives
\[
|\noderef{TwoBiteX}(\omega_w,I,x_i)|
\le \noderef{TwoBiteSmallCutoff}(\varepsilon,n)=n^{2\varepsilon}.
\]
The cardinality bound \(\noderef{FinalPairsCardLeHalfSq}\) gives
\[
|F_{x_i}(w)|
\le
|\noderef{TwoBiteFinalPairs}(\noderef{TwoBiteX}(\omega_w,I,x_i))|
\le \frac12\,|\noderef{TwoBiteX}(\omega_w,I,x_i)|^2
\le \frac12 n^{4\varepsilon}=c_{\mathrm{prod}}.
\]
This applies to \(w=v\) and to \(w=v'\).  Hence both terms in the maximum are
at most \(c_{\mathrm{prod}}\), and the absolute difference of the two counted
values is at most \(c_{\mathrm{prod}}\).
\end{proof}
